\scanpage{47}

Re NB \qquad Considérons

\[
\deleted{\unsure{f}}\varphi_N:\mathscr A\longrightarrow\mathscr A^\vee=E\qquad\text{via forme bil. ass. à forme qu. }N
\]

\[
\text{d'où}\qquad \mathscr A\xrightarrow{\deleted{\unsure{f}}\varphi_N}E\xrightarrow{Q}L^{\otimes 2}
\]

On trouve

\[
\boxed{Q(\varphi_N(x))=\deleted{\unsure{N(x)}}N(x)\delta}
\]

\margin{Vérification à partir du cas décomposé,}

\deleted{Si $\delta$ inv. (cas étale)}

Conditions équivalentes

a) $X'$ rev. \uline{étale} de $X$

b) $\varphi_N:\mathscr A\longrightarrow\mathscr A^\vee=E$ \uline{iso}

c) $\delta\in\Gamma L^{\otimes 2}$ section \uline{inv.} i.e. ne s'annule en aucun pt

Alors $\delta$ base de $L^{\otimes 2}$, $\varphi_N:\mathscr A\simeq E$, et en faisant les identifications, $Q$ correspond à $-N$

Si par contre $\delta=0$ (cas \deleted{\unsure{tott|tot\textsuperscript{t}}} ramifié) on trouve

\[
Q(\varphi_N(x))=0\ ..
\]

Considérons

\[
x\longmapsto\boxed{2x-T=z}\quad\inserted{(z\in\Gamma L\deleted{\unsure{^{\otimes 2}}})}\qquad \underline P_1\longrightarrow\underline L
\]

c'est un iso si 2 inv. sur $X$ donc iso sur $X_2=X-V(2)$. Si $x$ et $z$ se correspondent, on a

\[
2x=z+T\quad\Longrightarrow\quad Q(z+T)=4Q(x)
\]

et \deleted{le} \inserted{si $x$ section du} rev. quadratique $X'$ de $X$ i.e. \unsure{$\sqrt[2]{Q(x)=0}$|$Q(x)=0$}, on a $Q(T+z)=0$. Or on trouve \deleted{pour} \inserted{(comme $\psi(z+T)=2$)}

\[
\begin{aligned}
&\deleted{Q(T+z)}\\
&Q(z+T)=z^2\deleted{\unsure{+}} \unsure{-}\delta\\
&z^2+2Tz+\deleted{\unsure{T^2}}\unsure{-2}(Tz\unsure{-|+}\deleted{\unsure{T^2}})+4N=z^2\unsure{-}(T^2-4N)=z^2-\delta
\end{aligned}
\]

donc

\[
\left\{\begin{aligned}
&Q(x)=0\quad\text{i.e. }x\in\Gamma(X'/X)\quad\Longrightarrow\quad\boxed{z^2=\delta}\\
&(\text{et inv. si }2\text{ inv. sur }X.)
\end{aligned}\right.
\]

% archive sheet number: 46
